\documentclass[11pt]{article}
\usepackage[T1]{fontenc}
\usepackage[utf8]{inputenc}
\usepackage{lmodern,amsmath,amssymb,amsthm,mathtools}
\usepackage[margin=1in]{geometry}
\usepackage[hidelinks]{hyperref}
\usepackage{microtype}
\newtheorem{theorem}{Theorem}[section]
\newtheorem{lemma}[theorem]{Lemma}
\newtheorem{proposition}[theorem]{Proposition}
\newtheorem{corollary}[theorem]{Corollary}
\theoremstyle{remark}
\newtheorem{remark}[theorem]{Remark}
\newcommand{\dd}{\,d}
\newcommand{\Ran}{\operatorname{Ran}}
\newcommand{\supp}{\operatorname{supp}}
\newcommand{\tr}{\operatorname{tr}}
\newcommand{\sgn}{\operatorname{sgn}}
\newcommand{\R}{\mathbb R}
\newcommand{\C}{\mathbb C}
\title{Sharp logarithmic coupling estimates for theta equality frames\\
in a signed Weil-form realization}
\author{Leonard van Hemert}
\date{11 October 2026}
\begin{document}
\maketitle
\begin{abstract}
We determine the size of the finite-frame Schur correction associated
with the odd theta derivatives in a weighted realization of Weil's
quadratic form. For the first \(m\) pole-killing odd derivatives,
the exact correction satisfies
\[
 \|E_{m,\beta}\|\asymp_m
       \beta\bigl(\log(1/\beta)\bigr)^{4m+5/2}
       \qquad(\beta\downarrow0).
\]
Both inequalities are unconditional, with explicit constants and
parameter regimes. The upper bound follows from a constructive
cancellation of the source charge on each half-line, a square-root
range argument, and an integrated estimate of the full theta tail.
Odd box tests give the matching lower bound. We also obtain a
pointwise upper bound for the low-energy spectral mass of the
weighted gamma operator. The prime-shift operator and the polar
terms are retained in the exact Schur reduction. The estimate rules
out an \(O_m(\beta)\) bound for this correction; it does not determine
the sign of the remaining complementary form. An exact change of
metric also identifies how eliminating the equality space before
gamma normalization cancels the explicit correction and moves
its logarithmic growth into the metric.
\end{abstract}

\section{The estimate and its setting}

We use
\[
 \xi(s)=\tfrac12s(s-1)\pi^{-s/2}\Gamma(s/2)\zeta(s),\qquad
 \widehat f(t)=\int_{\R}f(x)e^{itx}\dd x.
\]
The Fourier convention has Plancherel factor \(1/(2\pi)\).
Inner products are linear in the first entry. The actual even
theta kernel is
\begin{equation}\label{phi}
 \Phi(x)=\sum_{n\ge1}
 (4\pi^2n^4e^{9|x|/2}-6\pi n^2e^{5|x|/2})
                    e^{-\pi n^2e^{2|x|}},
 \qquad \widehat\Phi(t)=\xi(1/2+it).
\end{equation}
It is smooth and positive. Its smooth even extension and Fourier
normalization are classical; see Romik
\cite[(1.8)--(1.11), printed p.~2]{Romik}.
Put
\[
 c(x)=\cosh(x/2),\qquad w(x)^2=\Phi(x)/c(x),\qquad H=L^2(\R,\dd x).
\]
Let \(\psi=\Gamma'/\Gamma\), and define the complete gamma symbols
\begin{align}
 m_\Gamma(t)&=\Re\psi(1/4+it/2)-\log\pi,\notag\\
 c_\Gamma&=m_\Gamma(0)
   =-\gamma_E-\pi/2-3\log2-\log\pi,\notag\\
 b(t)&=m_\Gamma(t)-c_\Gamma
 =\sum_{j\ge0}\frac{(t/2)^2}
 {(j+1/4)\{(j+1/4)^2+(t/2)^2\}}.\label{symbol}
\end{align}
Define \(G\) through its maximal closed nonnegative form
\begin{equation}\label{g}
 g(u)=\frac1{2\pi}\int_\R b(t)|\widehat{wu}(t)|^2\dd t,\qquad
 \mathcal V=\{u\in H:g(u)<\infty\}.
\end{equation}
The form domain, injectivity and relevant operator domains are proved
below. No critical-line location of a zero is assumed.

For \(r\ge1\) set
\begin{equation}\label{jets}
 d_r=2r+1,\qquad \kappa_r=4r+5/2,\qquad
 h_r=\partial_x(\partial_x^2-1/4)^r\Phi,\qquad
 \widetilde u_r=h_r/w.
\end{equation}
These are linearly independent odd vectors in \(D(G)\).
Fix \(m\ge1\), and choose any base-\(L^2\) orthonormal basis
\(u_1,\ldots,u_m\) of their first-\(m\) span. Write
\(U:\C^m\to H_{\rm odd}\) for the associated isometry, and put
\begin{align}
 F_\beta&=(G+\beta)^{-1/2},&
 B_\beta&=(G+\beta)^{1/2}U,\notag\\
 H_0&=U^*GU,& H_\beta&=H_0+\beta I,\notag\\
 W_\beta&=B_\beta H_\beta^{-1/2},&
 N_\beta&=I-W_\beta W_\beta^*.\label{frame}
\end{align}
The finite matrix \(H_0\) is positive definite. Thus \(W_\beta\)
is an isometry, and \(N_\beta\) is an orthogonal projection.
The finite-rank correction of interest is
\begin{equation}\label{E}
 E_{m,\beta}=\beta N_\beta F_\beta UU^*F_\beta N_\beta
                 \quad\text{on }H_{\rm odd}.
\end{equation}

\begin{theorem}[Sharp fixed-frame coupling]\label{main}
For every fixed \(m\ge1\) there are explicitly defined positive
constants \(a_m,A_m,\ell_m\) such that
\begin{equation}\label{mainbound}
 a_m\beta\bigl(\log(1/\beta)\bigr)^{\kappa_m}
 \le \|E_{m,\beta}\|
 \le A_m\beta\bigl(\log(1/\beta)\bigr)^{\kappa_m}
       \quad(0<\beta\le e^{-\ell_m}).
\end{equation}
One may take the constants in \eqref{frameconstants}--\eqref{regime}.
In the same regime,
\begin{equation}\label{mass}
 \|U^*\mathbf1_{(0,t]}(G)U\|
      \le 2A_m t\bigl(\log(1/t)\bigr)^{\kappa_m}
           \quad(0<t\le e^{-\ell_m}).
\end{equation}
The exact norm identity is
\begin{equation}\label{Enorm}
 \|E_{m,\beta}\|
   =\beta\|U^*(G+\beta)^{-1}U-H_\beta^{-1}\|.
\end{equation}
\end{theorem}

In particular, the first vector gives the exponent \(13/2\).
The constants and the small-parameter regime depend on \(m\).
Equation \eqref{mass} is an upper bound for a discontinuous spectral
projection; a matching pointwise lower bound is not asserted.

The estimate concerns a specific full-symbol resolvent and its exact
finite-dimensional correction. The Weil and Li frameworks themselves
have substantial antecedents. Bombieri--Lagarias
\cite[Theorem~1 and (2.4); Section~4]{BL} give the explicit-formula
relation underlying Li's criterion. Suzuki's norm formula
\cite[Theorem~1.1, (1.8)]{SuzukiLi} is equivalent to RH, and its
orthonormal-basis argument
\cite[Proposition~3.2]{SuzukiLi} assumes RH. Suzuki's localized
Weil operators \cite[Theorem~1.1; Sections~2.4--2.5]{SuzukiWeil}
provide the local gamma and prime formulas used here. Our weighted
whole-line \(G\) is instead defined directly by the closed form
\eqref{g}. The proof below quantifies its actual theta-jet resolvent,
then identifies the correction in the full signed Weil form.
No inner-function or conditional model-space property is imported.
Invertible-block Schur elimination is standard; compare
Dusson--Sigal--Stamm \cite[Theorem~1.2, (1.10)--(1.11)]{Feshbach}.
The quantitative step here is the actual theta-weighted resolvent
estimate. The zero-charge primitive is a Fourier antiderivative
construction, used with separate half-line compensation to obtain
the stated rate.

\section{The complete weighted gamma form}

\begin{lemma}[Theta bounds]\label{weights}
With \(C_-=4\pi^2-6\pi>0\) and \(C_\Phi=8\pi^2\), for every real \(x\),
\begin{align}
 C_-e^{9|x|/2-\pi e^{2|x|}}
 &\le \Phi(x)\le C_\Phi e^{9|x|/2-\pi e^{2|x|}},\label{phibound}\\
 C_-e^{4|x|-\pi e^{2|x|}}
 &\le w(x)^2\le2C_\Phi e^{4|x|-\pi e^{2|x|}}.\label{wbound}
\end{align}
For every fixed \(j\ge0\), a finite constant \(C_j\) gives
\[
 |\Phi^{(j)}(x)|\le C_j e^{(9/2+2j)|x|-\pi e^{2|x|}}.
\]
\end{lemma}
\begin{proof}
For \(x\ge0\), every theta summand is positive. Its \(n=1\) term
gives the lower bound. For the upper bound discard the negative
terms and factor out \(e^{9x/2-\pi e^{2x}}\).
The remaining sum is
\(4\pi^2\sum n^4e^{-\pi(n^2-1)e^{2x}}\).
For \(n\ge2\), \(n^4\le16\,6^{n-2}\) and \(n^2-1\ge n+1\).
Since \(q=e^{-\pi}<1/20\), this tail is at most
\(16q^3/(1-6q)<1\). The total sum is less than two.
Divide by \(c(x)\), using \(e^{x/2}/2\le c(x)\le e^{x/2}\).
Evenness supplies the other half-line.

Each fixed derivative of a summand is a polynomial of degree at most
\(j\) in \(\pi n^2e^{2x}\), times the original exponential.
Factoring the first exponential leaves a summable Gaussian
\(n\)-polynomial majorant. This proves the derivative bound and
justifies termwise differentiation locally. In particular every
fixed derivative is Schwartz.
\end{proof}

\begin{lemma}[Closed form, jump formula and domains]\label{closed}
The form \eqref{g} is densely defined and closed, and its associated
self-adjoint operator \(G\) is nonnegative and injective.
For \(f=wu\),
\begin{equation}\label{jump}
 g(u)=\frac12\int_\R
       \frac{e^{-|a|/2}}{1-e^{-2|a|}}
          \int_\R |f(x+a)-f(x)|^2\dd x\,\dd a.
\end{equation}
Compact smooth functions are a form core.
If \(u\in\mathcal V\) and the physical Fourier multiplier
\(b(D_x)(wu)\) belongs to \(L^2\), then
\[
 u\in D(G),\qquad Gu=w\,b(D_x)(wu).
 \tag{\(\ast\)}
\]
Moreover, at every \(t\ne0\),
\begin{equation}\label{inverse}
 b(t)\ge\frac{16t^2}{1+4t^2},\qquad
 b(t)^{-1}\le\frac14+\frac1{16t^2}.
\end{equation}
\end{lemma}
\begin{proof}
The digamma series gives \eqref{symbol}; its first nonnegative
summand gives \eqref{inverse}. Also
\[
 b(t)=\int_\R \mathfrak h(a)(1-\cos ta)\dd a,\qquad
 \mathfrak h(a)=\frac{e^{-|a|/2}}{1-e^{-2|a|}}.
\]
Indeed expand \(\mathfrak h=\sum_{j\ge0}e^{-(2j+1/2)|a|}\),
integrate each nonnegative term, and obtain \eqref{symbol}.
Tonelli and Plancherel prove \eqref{jump}.
The series also shows \(b(t)=O(\log(2+|t|))\):
split it at \(j+1/4=|t|/2\), bounding the first part by a
harmonic sum and the second by a cubic tail.

Multiplication by the bounded \(w\) maps \(H\) continuously into
\(L^2\). The map \(u\mapsto b(D_x)^{1/2}(wu)\) is closed:
convergence of \(u\) in \(H\) implies convergence of \(wu\), and
the physical multiplier is closed. This proves closedness of
\eqref{g}. Compact smooth functions give density.
If \(g(u)=0\), then \(\widehat{wu}\) is supported on \(\{0\}\),
because \(b(t)>0\) off zero. An \(L^2\) function with that support
vanishes. Since \(w>0\), \(u=0\).

Here are details of the core assertion. Let \(\chi_R\) be smooth,
equal to one on \([-R,R]\), supported on \([-2R,2R]\), and with
\(|\chi_R'|\le C/R\), for \(R\ge1\).
In the jump formula for \((1-\chi_R)f\), split its difference into
a cutoff times the difference of \(f\), and \(f\) times a cutoff
difference. Dominated convergence handles the first term.
The second tends to zero and is dominated by
\(C'\|f\|_2^2\mathfrak h(a)\min(a^2,1)\), an integrable function.
Thus \(\chi_Ru\to u\) in the form norm.
For a fixed \(R\), mollify the compact physical function
\(f_R=w\chi_Ru\). Fourier dominated convergence gives convergence
of the mollification in the \(b\)-form norm as well as in \(L^2\).
Divide the mollified functions by \(w\); on the fixed enlarged
compact support, \(1/w\) is smooth and bounded. This gives compact
smooth approximants to \(\chi_Ru\) in the required form norm.
Finally, if \(b(D_x)(wu)\in L^2\), Plancherel gives the polarized
form as \(\langle w b(D_x)(wu),v\rangle\) for every
\(v\in\mathcal V\). The form representation theorem proves
\((\ast)\).
\end{proof}

\section{Charge cancellation and the resolvent upper bound}

\begin{lemma}[Constructive square-root range approximation]\label{charge}
Suppose \(u\in H\), \(a=u/w\) is odd, and
\(|a(x)|\le Ce^{\kappa|x|}\), where \(C,\kappa>0\).
Complex-valued odd sources are allowed.
Put
\[
 R_\kappa=\max\left(1,\frac12\log\frac{2(\kappa+1)}{\pi}\right).
\]
For \(R\ge R_\kappa\), define
\begin{align}
 \alpha_R&=\int_0^R a(x)\dd x,&
 b_R&=\mathbf1_{[R+1,R+2]}-\mathbf1_{[-R-2,-R-1]},\notag\\
 a_R&=a\mathbf1_{[-R,R]}-\alpha_R b_R,&
 u_R&=wa_R,\qquad P_R(x)=\int_0^x a_R(t)\dd t.\label{charges}
\end{align}
Both half-line integrals of \(a_R\) vanish. The vector \(u_R\)
belongs to \(\Ran G^{1/2}\), and
\begin{align}
 \|G^{-1/2}u_R\|^2&\le K_Ae^{2\kappa R},\label{rangebound}\\
 \|u-u_R\|^2&\le K_Te^{(2\kappa+2)R-\pi e^{2R}},\label{tailbound}\\
 K_A&=C^2\left(\frac1{4\kappa}+\frac3{4\kappa^2}
                                +\frac1{16\kappa^3}\right),\notag\\
 K_T&=8C_\Phi C^2\left(\frac1\pi+\frac1{\kappa^2}\right).\notag
\end{align}
\end{lemma}
\begin{proof}
Oddness cancels the negative half-line charge once the positive one
is canceled. The primitive \(P_R\) is even and supported in
\([-R-2,R+2]\). On \([0,R]\),
\(|P_R(x)|\le Ce^{\kappa x}/\kappa\); on \([R,R+1]\) it is
\(\alpha_R\); on \([R+1,R+2]\) it decreases linearly to zero.
The supports of the two pieces of \(a_R\) are disjoint. Therefore
\begin{align*}
 \|a_R\|^2&\le C^2(\kappa^{-1}+2\kappa^{-2})e^{2\kappa R},\\
 \|P_R\|^2&\le C^2(\kappa^{-3}+4\kappa^{-2})e^{2\kappa R}.
\end{align*}
The derivative \(P_R'=a_R\) holds in distributions and almost
everywhere. Hence \(\widehat a_R=-it\widehat P_R\), and
\eqref{inverse} gives
\[
 \|b(D_x)^{-1/2}a_R\|^2
 \le\tfrac14\|a_R\|^2+\tfrac1{16}\|P_R\|^2
 \le K_Ae^{2\kappa R}.
\]
For every \(v\in\mathcal V\), Fourier Cauchy--Schwarz yields
\[
 |\langle u_R,v\rangle|
 =|\langle a_R,wv\rangle|
 \le\|b(D_x)^{-1/2}a_R\|\,\|G^{1/2}v\|.
\]
Injectivity makes \(\Ran G^{1/2}\) dense. The bounded functional
on this range has a Riesz representative \(q_R\) of norm at most
\(\|b(D_x)^{-1/2}a_R\|\). Thus
\(\langle u_R,v\rangle=\langle q_R,G^{1/2}v\rangle\)
on the entire form domain. The adjoint-domain characterization of
the self-adjoint \(G^{1/2}\) gives
\(q_R\in D(G^{1/2})\) and \(G^{1/2}q_R=u_R\).
This proves \eqref{rangebound} without taking the formal inverse
of a product of unbounded operators.

Write \(Y=e^{2R}\). By \eqref{wbound}, the omitted source tail has
weighted square norm at most
\[
 4C_\Phi C^2\int_R^\infty
 e^{(2\kappa+4)x-\pi e^{2x}}\dd x
 =2C_\Phi C^2\int_Y^\infty y^{\kappa+1}e^{-\pi y}\dd y
 \le\frac{4C_\Phi C^2}{\pi}Y^{\kappa+1}e^{-\pi Y}.
\]
Here \(\pi Y\ge2(\kappa+1)\), so the logarithmic derivative of
the last integrand is at most \(-\pi/2\).
The weighted square norm of the outer charge rectangles is at most
\[
 \frac{4C_\Phi C^2}{\kappa^2}
       e^{(2\kappa+4)R+4-\pi e^2Y}
 \le\frac{4C_\Phi C^2}{\kappa^2}Y^{\kappa+1}e^{-\pi Y}.
\]
For the final inequality use
\(Ye^4e^{-\pi(e^2-1)Y}\le1\) for \(Y\ge1\);
\(\log Y\le Y\), \(e^2>7\), and \(\pi>3\) suffice.
The square triangle inequality proves \eqref{tailbound}.
The half-line cancellation keeps the primitive's large plateau in
two outer strips of bounded length, instead of across the central
interval. This is why its norm has no additional factor \(R\).
\end{proof}

\begin{proposition}[Explicit upper rate]\label{upper}
Under the hypotheses of Lemma~\ref{charge}, define
\begin{equation}\label{KU}
 \ell_\kappa=\max(\pi e^2/2,\kappa+1),\qquad
 K_U(C,\kappa)=2(K_A+K_T)(2/\pi)^\kappa.
\end{equation}
For \(0<\beta\le e^{-\ell_\kappa}\),
\[
 \mathcal R_u(\beta):=\langle u,(G+\beta)^{-1}u\rangle
       \le K_U(C,\kappa)\bigl(\log(1/\beta)\bigr)^\kappa.
\]
Also, for \(0<t\le e^{-\ell_\kappa}\),
\[
 \langle\mathbf1_{(0,t]}(G)u,u\rangle
       \le2K_U(C,\kappa)t\bigl(\log(1/t)\bigr)^\kappa.
\]
\end{proposition}
\begin{proof}
The resolvent square-root triangle inequality gives
\[
 \mathcal R_u(\beta)\le
 2K_Ae^{2\kappa R}
       +2\beta^{-1}K_Te^{(2\kappa+2)R-\pi e^{2R}}.
\]
Set \(L=\log(1/\beta)\) and \(e^{2R}=2L/\pi\).
The stated regime is exactly \(R\ge R_\kappa\).
The error term is \(2K_T\beta(2L/\pi)^{\kappa+1}\).
Since \(\beta(2L/\pi)<1\), \eqref{KU} gives the first bound.
For the second apply
\(\mathbf1_{(0,t]}(s)\le2t/(s+t)\) in the spectral theorem.
No assumption on the type of the spectral measure is used.
\end{proof}

\section{Actual theta jets and matching lower bounds}

The following finite definitions supply constants for every jet.
For \(\alpha=9/2,5/2\), let
\begin{align}
 P_{\alpha,0}(q)&=1,&
 P_{\alpha,k+1}(q)&=(\alpha-2q)P_{\alpha,k}(q)
                                      +2qP'_{\alpha,k}(q),\notag\\
 R_{\alpha,r}(q)&=\sum_{k=0}^r\binom rk(-1/4)^{r-k}
                                  P_{\alpha,2k+1}(q),&
 B_{\alpha,r}&=\sum_{j=0}^{d_r}
          |[q^j]R_{\alpha,r}(q)|\,\pi^{j-d_r}.\label{polynomials}
\end{align}
For a positive integer \(p\), set
\begin{equation}\label{Sp}
 M=p+1,\qquad \rho_p=e^{1-\pi(2M+1)},\qquad
 \mathcal S_p=\sum_{n=1}^{M-1}n^pe^{-\pi(n^2-1)}
      +\frac{M^pe^{-\pi(M^2-1)}}{1-\rho_p}.
\end{equation}
Finally put
\begin{equation}\label{Ca}
 C_{a,r}=\frac{\pi^{d_r}}{C_-}
 \left(4\pi^2B_{9/2,r}\mathcal S_{4+2d_r}
              +6\pi B_{5/2,r}\mathcal S_{2+2d_r}\right).
\end{equation}

\begin{lemma}[Whole-lattice jet bounds]\label{jetbounds}
The constants \eqref{polynomials}--\eqref{Ca} give
\begin{equation}\label{sourceupper}
 \left|\frac{h_r(x)}{\Phi(x)}\right|
        \le C_{a,r}e^{2d_r|x|},\qquad
 \left|\frac{\widetilde u_r(x)}{w(x)}\right|
        \le C_{a,r}e^{\kappa_r|x|}.
\end{equation}
Every \(\widetilde u_r\) belongs to \(D(G)\).
Moreover, as \(x\to+\infty\),
\begin{equation}\label{jetleading}
 \frac{c(x)h_r(x)}{\Phi(x)}
 =-\frac{(2\pi)^{d_r}}2 e^{\kappa_rx}
                  (1+e^{-x})(1+O_r(e^{-2x})).
\end{equation}
Consequently the vectors \(\widetilde u_1,\ldots,\widetilde u_m\)
are independent for every \(m\).
\end{lemma}
\begin{proof}
The recurrence in \eqref{polynomials} is differentiation of
\(e^{\alpha x-\pi n^2e^{2x}}\).
The polynomial \(R_{\alpha,r}\) has degree \(d_r\) and leading
coefficient \(-2^{d_r}\). For \(q\ge\pi\),
\(|R_{\alpha,r}(q)|\le B_{\alpha,r}q^{d_r}\).
The \(n\)th differentiated theta term is
\[
 4\pi^2n^4e^{9x/2-q_n}R_{9/2,r}(q_n)
       -6\pi n^2e^{5x/2-q_n}R_{5/2,r}(q_n),
       \qquad q_n=\pi n^2e^{2x}.
\]
For \(n\ge M=p+1\), the ratio of successive terms of
\(n^pe^{-\pi(n^2-1)}\) is at most
\(e^{p/n-\pi(2n+1)}\le\rho_p<1\).
Thus \(\mathcal S_p\) bounds the complete Gaussian sum.
Divide the derivative series by the positive lower bound
\eqref{phibound}, and use
\(e^{-\pi(n^2-1)e^{2x}}\le e^{-\pi(n^2-1)}\) for \(x\ge0\).
This proves the first inequality in \eqref{sourceupper}.
Multiplication by \(c(x)\le e^{|x|/2}\) proves the second.
Parity supplies \(x<0\).

The lower weight bound and the derivative bounds show
\(\widetilde u_r=h_r/w\in L^2\).
The physical function \(h_r\) is Schwartz, and
\(b(t)=O(\log(2+|t|))\), so \(b(D_x)h_r\in L^2\).
Lemma~\ref{closed} puts \(\widetilde u_r\) in \(D(G)\).

In the \(n=1\) term, take the leading coefficients of the finite
polynomials and divide by the leading term of \(\Phi\).
The relative polynomial error is \(O_r(e^{-2x})\).
For \(n\ge2\), the complete relative remainder is bounded by a
constant times
\[
 \sum_{n\ge2}n^{4+2d_r}e^{-\pi(n^2-1)e^{2x}}
 \le \mathcal S_{4+2d_r}
                e^{-3\pi(e^{2x}-1)}.
\]
It is smaller than the stated polynomial error. Multiplication by
\(c(x)=e^{x/2}(1+e^{-x})/2\) proves \eqref{jetleading}
with the full lattice retained.
A nontrivial finite linear combination has a unique largest
\(\kappa_r\); its leading tail cannot vanish. This proves independence.
\end{proof}

To specify the lower regime, define
\begin{align}
 S_{\alpha,r}&=\sum_{j<d_r}|[q^j]R_{\alpha,r}(q)|,\notag\\
 Q_r&=\max(1,2^{1-d_r}S_{9/2,r},2^{1-d_r}S_{5/2,r}),\notag\\
 R_{0,r}&=\max(1,\tfrac12\log(Q_r/\pi)),&
 c_{a,r}&=2^{d_r-4}\pi^{d_r},\notag\\
 R_{*,r}&=R_{0,r}+\kappa_r^{-1}
                    \log\frac{2(C_{a,r}+c_{a,r})}{c_{a,r}},\notag\\
 \overline C_g&=218/25,&
 \ell_{L,r}&=2\pi e^{2R_{*,r}},\notag\\
 K_{L,r}&=\frac{c_{a,r}^2}
 {\kappa_r^2(4\overline C_g+1)(2\pi)^{\kappa_r}}.\label{lowerconstants}
\end{align}

\begin{proposition}[Explicit matching lower rate]\label{lower}
For every \(r\ge1\) and \(0<\beta\le e^{-\ell_{L,r}}\),
\[
 \mathcal R_{\widetilde u_r}(\beta)
       \ge K_{L,r}\bigl(\log(1/\beta)\bigr)^{\kappa_r}.
\]
\end{proposition}
\begin{proof}
For \(q\ge Q_r\), \(q\ge1\) and the sum of all lower-degree
terms is at most \(S_{\alpha,r}q^{d_r-1}\).
Therefore
\[
 -3\cdot2^{d_r-1}q^{d_r}
       \le R_{\alpha,r}(q)\le-2^{d_r-1}q^{d_r}.
\]
For \(x\ge R_{0,r}\), this holds for every \(q_n\).
In each differentiated theta summand, the ratio of the possible
positive contribution from the second term to the negative
leading bound from the first is at most
\[
 \frac{9}{2\pi n^2}e^{-2x}<\tfrac12.
\]
Thus all derivative summands are negative. Retaining half of the
first negative bound and then the \(n=1\) term yields
\[
 h_r(x)\le-2^{d_r}\pi^{d_r+2}
                  e^{(9/2+2d_r)x-\pi e^{2x}}.
\]
Divide this negative numerator by the \emph{upper} bound on
\(\Phi\), and use \(c(x)\ge e^{x/2}/2\). With
\(a_r=c h_r/\Phi=\widetilde u_r/w\), we obtain
\begin{equation}\label{negative}
                 a_r(x)\le-c_{a,r}e^{\kappa_rx}
                             \quad(x\ge R_{0,r}).
\end{equation}

Let \(H_R(x)=\sgn(x)\mathbf1_{[-R,R]}(x)\) and \(v_R=H_R/w\).
For finite \(R\), \(v_R\in L^2\).
The physical box \(H_R\) has supremum one and total variation four.
Hence
\[
 \int|H_R(x+a)-H_R(x)|^2\dd x\le8|a|.
\]
The full jump formula gives
\[
 g(v_R)\le4 C_g,\qquad
 C_g=\int_\R\mathfrak h(a)|a|\dd a
     =2\sum_{j\ge0}(2j+1/2)^{-2}
     \le8+2(4/25+1/5)=\overline C_g.
\]
This also proves \(v_R\in\mathcal V\), including its jumps.
The upper central bound \eqref{sourceupper} and the negative tail
\eqref{negative} imply
\[
 \langle\widetilde u_r,v_R\rangle
 =2\int_0^R a_r(x)\dd x
 \le-\frac{c_{a,r}}{\kappa_r}e^{\kappa_rR}
                        \quad(R\ge R_{*,r}).
\]
Indeed the central integral contributes at most
\[
 \frac{2C_{a,r}}{\kappa_r}e^{\kappa_rR_{0,r}},
\]
while the tail contributes at most
\[
 -\frac{2c_{a,r}}{\kappa_r}
       \bigl(e^{\kappa_rR}-e^{\kappa_rR_{0,r}}\bigr).
\]
The definition of \(R_{*,r}\) gives the displayed inequality.

The lower weight bound, and monotonicity of
\(e^{-4x+\pi e^{2x}}\) on \(x\ge0\), give
\[
 \|v_R\|^2\le\frac{2R}{C_-}e^{-4R+\pi e^{2R}}.
\]
Set \(L=\log(1/\beta)\) and \(e^{2R}=L/(2\pi)\).
Then \(R\ge R_{*,r}\ge1\) in the stated regime, and
\[
 \beta\|v_R\|^2
 \le\frac{2R}{C_-}e^{-4R-L/2}
 \le\frac{2e^{-4}}{C_-}<1.
\]
The resolvent variational inequality on the form domain now gives
\[
 \mathcal R_{\widetilde u_r}(\beta)
 \ge
 \frac{|\langle\widetilde u_r,v_R\rangle|^2}
 {g(v_R)+\beta\|v_R\|^2}
 \ge
 \frac{c_{a,r}^2}{\kappa_r^2(4\overline C_g+1)}
                     \left(\frac{L}{2\pi}\right)^{\kappa_r}.
\]
It follows directly by Cauchy--Schwarz applied to
\((G+\beta)^{-1/2}\widetilde u_r\) and
\((G+\beta)^{1/2}v_R\).
This proves the proposition with no truncation of the gamma kernel.
\end{proof}

\section{The exact signed Weil coupling}

This section identifies the correction \eqref{E} in the complete
Weil form, including the two polar moments. It is not needed for
the scalar resolvent estimates.

Define the bounded self-adjoint full prime-shift operator by
\begin{equation}\label{prime}
 (Ju)(x)=\sum_{n\ge2}\frac{\Lambda(n)}{\sqrt n}w(x)
   \{w(x-\log n)u(x-\log n)+w(x+\log n)u(x+\log n)\}.
\end{equation}
For \(q\in H\), let \((q\otimes q)u=\langle u,q\rangle q\).
Put
\begin{align}
 p_c&=w\cosh(x/2),&p_s&=w\sinh(x/2),\notag\\
 \mathcal P&=2(p_c\otimes p_c-p_s\otimes p_s),&
 D&=J-c_\Gamma M_w^2,\qquad D_W=D-\mathcal P.\label{DW}
\end{align}
Both pole vectors belong to \(H\).

\begin{lemma}[Full operator and unconditional equality]\label{equality}
The series \eqref{prime} converges in operator norm, and \(D_W\)
is bounded and self-adjoint. On compact smooth physical tests \(f=wu\),
the Weil form is exactly
\begin{equation}\label{weil}
 Q_W(f)=g(u)-\langle D_Wu,u\rangle.
\end{equation}
This formula supplies a closed lower-bounded weighted extension.
Every vector in \eqref{jets} satisfies, unconditionally,
\begin{equation}\label{annihilate}
 \mathcal P\widetilde u_r=0,\qquad
 (G-D)\widetilde u_r=(G-D_W)\widetilde u_r=0.
\end{equation}
\end{lemma}
\begin{proof}
For \(a=\log n\), set \(A=e^{2|x|}\), \(B=e^{2|x-a|}\).
By \eqref{wbound},
\[
 w(x)w(x-a)\le2C_\Phi AB e^{-\pi(A+B)/2}
                   \le C_\Phi e^{-\pi(A+B)/8}
                   \le C_\Phi e^{-\pi n/4}.
\]
Here \(A,B\ge1\), \(\log(2AB)\le A+B<3\pi(A+B)/8\),
and \(AB\ge n^2\), so \(A+B\ge2n\).
The paired shift has operator norm at most
\(2C_\Phi e^{-\pi n/4}\).
Since \(\Lambda(n)\le\log n\),
\(\sum_{n\ge2}2C_\Phi(\log n)n^{-1/2}e^{-\pi n/4}<\infty\).
The two shifts are adjoints of one another, proving the assertion.
The remaining terms in \eqref{DW} are bounded self-adjoint operators.

For completeness, specify the local explicit formula used here.
Let \(\mathcal Z\) be the distinct zeros, with multiplicities
\(m_\gamma\), of \(X(z)=\xi(1/2+iz)\).
They satisfy \(|\Im\gamma|<1/2\) without RH.
For \(b_0\in C_c^\infty(\R)\), Weil's explicit formula gives the
local distribution
\begin{equation}\label{distribution}
 \langle k,b_0\rangle
     =\sum_{\gamma\in\mathcal Z}m_\gamma\widehat b_0(\gamma),
 \qquad
 k=2\cosh(\,\cdot\,/2)+K_\Gamma
       -\sum_{n\ge2}\frac{\Lambda(n)}{\sqrt n}
                (\delta_{\log n}+\delta_{-\log n}),
\end{equation}
where \(K_\Gamma\) is the tempered Fourier multiplier
\(m_\Gamma(D_x)\), with its finite-part kernel at zero.
The zero sum is absolutely convergent for these tests.
This is the unconditional compact-test formula; see
Suzuki \cite[Theorem~1.1(2), (1.3)]{SuzukiScrew} and
\cite[Sections~2.4--2.5; (3.1), printed p.~13]{SuzukiWeil}.
Only \(K_\Gamma\), not the whole \(k\), is assigned a global
tempered Fourier multiplier.
One can obtain its linear zero formula directly from the cited
unconditional expansion
\[
 \Psi(t)=\sum_{\gamma\in\mathcal Z}m_\gamma
                   \frac{1-\cos(\gamma t)}{\gamma^2},
                       \qquad k=\Psi''.
\]
The series converges absolutely and uniformly on compact real
intervals, by the strip bound and
\(\sum m_\gamma(1+|\gamma|)^{-2}<\infty\).
There is no zero at \(0\), since \(X(0)=\int\Phi>0\).
Pair it with \(b_0''\), integrate twice by parts, and use the
symmetry of the full zero multiset under \(\gamma\mapsto-\gamma\).
This gives exactly \eqref{distribution}, without a half-set factor.
The contact normalization can also be specified directly:
\begin{align*}
 \langle K_\Gamma,b_0\rangle
 ={}&-(\log(4\pi)+\gamma_E)b_0(0)\\
 &-\int_0^\infty
   \{b_0(x)+b_0(-x)-2e^{-x/2}b_0(0)\}
               \frac{e^{-x/2}}{1-e^{-2x}}\dd x .
\end{align*}
The subtraction makes the integral convergent at zero.
The digamma integral confirms its Fourier symbol \(m_\Gamma\);
the constants agree because
\[
 2\int_0^\infty\frac{e^{-x}-e^{-2x}}{1-e^{-2x}}\dd x
       =2\log2=\log4.
\]
Suzuki's 2026 xi normalization omits the factor \(1/2\).
That changes neither the zeros, their multiplicities nor this
explicit formula; \eqref{phi} retains the normalization used here.

Apply \eqref{distribution} to the correlation
\(b_0(t)=\int f(x)\overline{f(x-t)}\dd x\).
Equivalently, for a compact smooth \(f\),
\[
 Q_W(f)=\sum_{\gamma\in\mathcal Z}m_\gamma
              \widehat f(\gamma)\overline{\widehat f(\bar\gamma)}.
\]
This sum is absolutely convergent by compact-test Fourier decay,
the strip bound and zero counting. Multiplicities are retained.
The polar contribution is
\[
 2\left|\int f(x)\cosh(x/2)\dd x\right|^2
       -2\left|\int f(x)\sinh(x/2)\dd x\right|^2.
\]
The gamma contribution is
\((2\pi)^{-1}\int m_\Gamma(t)|\widehat f(t)|^2\dd t\);
the prime contribution is \(-\langle Ju,u\rangle\).
Splitting \(m_\Gamma=b+c_\Gamma\) proves \eqref{weil}.
Bounded perturbation of the closed form \(g\) gives its stated
weighted extension. The zero formula itself is only being invoked
on the compact tests.

Here is an extension of the local formula sufficient for equality.
For a fixed real \(x\), test with a compact smooth cutoff of
\(t\mapsto h_r(x-t)\).
For each derivative of order at most two, the cutoff remainder
tends to zero in \(L^1(e^{|t|/2}\dd t)\), by Lemma~\ref{weights}.
Two integrations by parts and the bound on \(|\Im\gamma|\) give
\[
 |\widehat{b_{x,R}}(\gamma)-\widehat{h_r(x-\cdot)}(\gamma)|
       \le \epsilon_{x,R}(1+|\gamma|)^{-2},\qquad
                                      \epsilon_{x,R}\to0.
\]
Zero counting gives
\(\sum m_\gamma(1+|\gamma|)^{-2}<\infty\).
The limiting transform is zero at every \(\gamma\), since
\[
 \widehat h_r(z)=(-iz)(-z^2-1/4)^rX(z).
\]
Evenness of \(X\) also makes it zero at \(-\gamma\).
Thus the zero side tends absolutely to zero.
On the physical side the polar term is integrable against \(h_r\);
the gamma cutoff converges in Schwartz seminorms; and the entire
prime series converges locally absolutely, with all fixed derivatives,
by its theta decay and \(\Lambda(n)\le\log n\).
Hence \(k*h_r=0\) locally on the whole real line.
No rearrangement into an unproved global Fourier transform of \(k\)
has been made.

Integration by parts gives
\[
 \int h_r(x)e^{\pm x/2}\dd x=0:
 \quad(\partial_x^2-1/4)^r e^{\pm x/2}=0
                      \quad(r\ge1).
\]
There are no boundary terms, by the full theta decay.
Both pole moments therefore vanish, and \eqref{distribution} gives
\[
 K_\Gamma h_r
   =\sum_{n\ge2}\frac{\Lambda(n)}{\sqrt n}
                \{h_r(\,\cdot-\log n)+h_r(\,\cdot+\log n)\}.
\]
Multiply by \(w\). Lemma~\ref{jetbounds} already puts
\(\widetilde u_r\) in \(D(G)\); the right side is the norm-convergent
operator \(J\widetilde u_r\).
Using \(K_\Gamma=b(D_x)+c_\Gamma\) proves \eqref{annihilate}.
\end{proof}

For this full signed operator, let
\begin{equation}\label{T}
 T_\beta=I-F_\beta D_WF_\beta\quad\text{on }H_{\rm odd}.
\end{equation}
All its factors are bounded for each fixed \(\beta>0\).
Reflection preserves the gamma form, the prime operator and the
pole operator, so the odd restriction is valid.

\begin{proposition}[Exact finite elimination]\label{schur}
The positive finite equality block is
\[
 W_\beta^*T_\beta W_\beta=\beta H_\beta^{-1}.
\]
Its Schur complement on \(N_\beta H_{\rm odd}\) is exactly
\begin{equation}\label{complement}
 \mathcal C_{m,\beta}
        =N_\beta T_\beta N_\beta-E_{m,\beta}.
\end{equation}
In particular their negative indices agree:
\(n_-(T_\beta)=n_-(\mathcal C_{m,\beta})\).
Identity \eqref{Enorm} holds, and its matrix difference is
nonnegative. The same correction results if \(D_W\) is replaced
by the gamma-prime operator \(D\).
\end{proposition}
\begin{proof}
Equality \eqref{annihilate} gives
\[
 T_\beta B_\beta=\beta F_\beta U,\qquad
 B_\beta^*F_\beta U=I.
\]
Thus \(W_\beta^*T_\beta W_\beta=\beta H_\beta^{-1}>0\).
In the orthogonal decomposition
\(H_{\rm odd}=\Ran W_\beta\oplus N_\beta H_{\rm odd}\),
complete the square with the inverse
\(\beta^{-1}H_\beta\) of this finite block.
The off-diagonal block is
\(\beta N_\beta F_\beta UH_\beta^{-1/2}\).
Its product with the inverse pivot is exactly \eqref{E};
the two factors \(H_\beta^{\pm1/2}\) cancel.
The completion is a bounded invertible triangular congruence
for each fixed \(\beta\), proving the negative-index assertion.

The nonzero eigenvalues of
\(N_\beta F_\beta UU^*F_\beta N_\beta\) equal those of
\[
 U^*F_\beta N_\beta F_\beta U
    =U^*(G+\beta)^{-1}U-H_\beta^{-1}\ge0.
\]
This proves \eqref{Enorm}. The calculation of this matrix needs
only the gamma form and its frame, rather than equality with a
signed operator. Equality is needed to identify it as the Schur
correction of \(T_\beta\).
Finally \(\mathcal P U=0\), so both signed choices satisfy
the same identity \(T_\beta B_\beta=\beta F_\beta U\).
The complementary block still contains every prime and polar term.
\end{proof}

\section{Proof of the fixed-frame rate and explicit constants}

Write \(u_j=\sum_{r=1}^m c_{jr}\widetilde u_r\), and define
\begin{align}
 C_j&=\sum_{r=1}^m|c_{jr}|C_{a,r},&
 A_m&=\sum_{j=1}^mK_U(C_j,\kappa_m),\notag\\
 q_m&=\|\widetilde u_m\|^2,&
 \mu_m&=\lambda_{\min}(H_0)>0,\notag\\
 a_m&=\frac{K_{L,m}}{2q_m}.\label{frameconstants}
\end{align}
One explicit regime is
\begin{equation}\label{regime}
 \ell_m=\max\left(\ell_{\kappa_m},\ell_{L,m},
                  \left(\frac{2q_m}{K_{L,m}\mu_m}\right)^{1/\kappa_m}
             \right).
\end{equation}
All definitions in these formulas are specified in
\eqref{polynomials}--\eqref{lowerconstants} and \eqref{KU}.

\begin{proof}[Proof of Theorem~\ref{main}]
Injectivity of \(G\) and independence of the finite frame imply
positive definiteness of \(H_0\). Since \(\kappa_r\le\kappa_m\),
\[
 |u_j(x)/w(x)|\le C_j e^{\kappa_m|x|}.
\]
Proposition~\ref{upper}, positivity of the resolvent matrix and
its trace yield
\[
 \|U^*(G+\beta)^{-1}U\|
     \le A_m\bigl(\log(1/\beta)\bigr)^{\kappa_m}
                         \quad(\log(1/\beta)\ge\ell_{\kappa_m}).
\]
The nonnegative matrix difference in \eqref{Enorm} is bounded
above by this matrix. This proves the upper estimate in
\eqref{mainbound}. The same trace and spectral inequality prove
\eqref{mass}.

The normalized vector \(\widetilde u_m/\sqrt{q_m}\) is in the
frame span. Proposition~\ref{lower} gives
\[
 \|U^*(G+\beta)^{-1}U\|
       \ge\frac{K_{L,m}}{q_m}
                      \bigl(\log(1/\beta)\bigr)^{\kappa_m}.
\]
Also \(\|H_\beta^{-1}\|\le\mu_m^{-1}\).
The reverse norm inequality, and the last term of \eqref{regime},
therefore give
\[
 \|U^*(G+\beta)^{-1}U-H_\beta^{-1}\|
       \ge\frac{K_{L,m}}{2q_m}
                      \bigl(\log(1/\beta)\bigr)^{\kappa_m}.
\]
Multiply by \(\beta\) and use \eqref{Enorm}.
Every bound holds on the stated full parameter interval.
\end{proof}

\begin{remark}[Evaluation of the frame constants]
The constants involve finite polynomials, finite sums, and
convergent integrals and matrices. They do not require a numerical
positivity assertion. For a definite choice of frame, let
\(V\) have columns \(\widetilde u_1,\ldots,\widetilde u_m\), set
\(M_0=V^*V\), and take \(U=VM_0^{-1/2}\).
Then \(c_{jr}=(M_0^{-1/2})_{rj}\) and
\[
 H_0=M_0^{-1/2}M_1M_0^{-1/2},\qquad
 M_1=V^*GV.
\]
The entries of \(M_0\) are ordinary theta integrals.
Those of \(M_1\) are the polarized full gamma integrals
\((2\pi)^{-1}\int b(t)\widehat h_r(t)
\overline{\widehat h_s(t)}\dd t\).
Their convergence was proved above.
If an expression without a smallest-eigenvalue symbol is desired,
use the explicit lower bound
\[
 \mu_m\ge \delta_m:=
       \frac{\det H_0}{(\tr H_0)^{m-1}}>0
\]
and replace \(\mu_m\) by \(\delta_m\) in \eqref{regime}.
The resulting larger threshold is valid for every \(m\), including
\(m=1\). These are effective fixed-frame definitions, with no
asserted uniform conditioning as \(m\) grows.
\end{remark}

\begin{remark}[The first jet]
The first polynomials are
\[
 R_{9/2,1}=90-183q+78q^2-8q^3,\qquad
 R_{5/2,1}=15-75q+54q^2-8q^3.
\]
One may use the general constant \eqref{Ca}, so no separate
numerical bound is required. The common formula gives
\(\kappa_1=13/2\), and all first-jet constants and its parameter
regime follow from the displayed finite definitions.
There is also a shorter lower-tail estimate for this jet.
For \(q\ge21\),
\begin{align*}
 &R_{9/2,1}(q)-\frac{3}{2q}R_{5/2,1}(q)
                       +4q^3\left(1-\frac3{2q}\right)\\
 &\hspace{2em}=-4q^3+84q^2-264q+\frac{405}{2}-\frac{45}{2q}<0.
\end{align*}
Since \(\pi e^2>21\), every theta summand at \(x\ge1\) has
derivative ratio at most \(-4(\pi n^2e^{2x})^3\).
All original summands are positive, so summing proves
\[
 \frac{\Phi'''(x)-\Phi'(x)/4}{\Phi(x)}
        \le-4\pi^3e^{6x},\qquad
 \frac{c(x)(\Phi'''(x)-\Phi'(x)/4)}{\Phi(x)}
        \le-2\pi^3e^{13x/2}\quad(x\ge1).
\]
Thus, for the first jet, the lower construction may instead use
\(R_{0,1}=1\) and \(c_{a,1}=2\pi^3\), leaving
\eqref{Ca} unchanged and recomputing
\(R_{*,1},\ell_{L,1},K_{L,1}\) by \eqref{lowerconstants}.
\end{remark}

\section{Normalization after equality-space elimination}

There is an exact way to place the large correction in context.
It eliminates the same equality space before gamma normalization.
The resulting signed form still contains the full arithmetic
operator, but the finite coupling term has disappeared.

\begin{proposition}[Exact change of metric]\label{metric}
Let \(P_0=UU^*\), \(N_0=I-P_0\) and \(Y=N_0H_{\rm odd}\).
Let \(G_\perp\) be the operator of the restricted gamma form on \(Y\),
and define
\[
 C=N_0GU,\qquad
 A_\perp=G_\perp-N_0D_WN_0.
\]
All occurrences of \(N_0D_WN_0\) below are restricted to \(Y\).
For \(\beta>0\) put
\begin{align}
 M_\beta&=G_\perp+\beta-CH_\beta^{-1}C^*\ge\beta,\notag\\
 J_\beta v&=v-UH_\beta^{-1}C^*v,\notag\\
 \mathcal O_\beta&=(G+\beta)^{1/2}J_\beta M_\beta^{-1/2}.
                                                        \label{shortmetric}
\end{align}
The map \(\mathcal O_\beta:Y\to N_\beta H_{\rm odd}\) is unitary,
and the exact complement obeys
\begin{equation}\label{metriccomplement}
 \mathcal O_\beta^*\mathcal C_{m,\beta}\mathcal O_\beta
       =M_\beta^{-1/2}(A_\perp+\beta)M_\beta^{-1/2}.
\end{equation}
The right side is understood by its bounded form.
With
\[
 \widehat T_\beta=(G_\perp+\beta)^{-1/2}
                    (A_\perp+\beta)(G_\perp+\beta)^{-1/2},
 \qquad
 V_\beta=(G_\perp+\beta)^{1/2}M_\beta^{-1/2},
\]
both operators are bounded, \(V_\beta\) is invertible, and the
right side of \eqref{metriccomplement} is
\(V_\beta^*\widehat T_\beta V_\beta\). Moreover
\begin{equation}\label{metricgrowth}
 \|V_\beta\|^2-1
 =\left\|H_\beta^{1/2}
       \{U^*(G+\beta)^{-1}U-H_\beta^{-1}\}H_\beta^{1/2}\right\|
 \asymp_m\bigl(\log(1/\beta)\bigr)^{\kappa_m}.
\end{equation}
\end{proposition}
\begin{proof}
The finite projection preserves the gamma form domain because its
columns lie in \(D(G)\). The restricted form on \(Y\) is closed
and densely defined. Its off-diagonal gamma block is the bounded
map \(C\), so the gamma block matrix is
\[
 G+\beta=
 \begin{pmatrix}H_\beta&C^*\\ C&G_\perp+\beta\end{pmatrix}
           \quad\text{on }\Ran U\oplus Y.
\]
Equality \((G-D_W)U=0\) and self-adjointness give
\[
 G-D_W=0\oplus A_\perp
\]
in the same decomposition. Completing the gamma square gives
\[
 \|(G+\beta)^{1/2}(Ua+v)\|^2
 =\|H_\beta^{1/2}(a+H_\beta^{-1}C^*v)\|^2
                         +\|M_\beta^{1/2}v\|^2.
\]
The infimum over \(a\) is at least \(\beta\|v\|^2\), since
\(G\ge0\) and \(\|Ua+v\|^2=\|a\|^2+\|v\|^2\).
Thus \(M_\beta\ge\beta\), and the lift in
\eqref{shortmetric} is an isometry into \(N_\beta H_{\rm odd}\).
It is onto: for \(z\in N_\beta H_{\rm odd}\),
\(x=F_\beta z=Ua+v\) lies in the form domain, and
\(B_\beta^*z=0\) gives
\(H_\beta a+C^*v=0\). Hence \(x=J_\beta v\).

Write \(x=J_\beta M_\beta^{-1/2}v\). The form of \(T_\beta\)
at \(\mathcal O_\beta v\) is the form of \(G-D_W+\beta\) at \(x\).
Writing \(\mathfrak a_\perp\) for the form of \(A_\perp\),
the exact kernel decomposition gives
\[
 \langle T_\beta\mathcal O_\beta v,\mathcal O_\beta v\rangle
 =\mathfrak a_\perp(M_\beta^{-1/2}v)
                   +\beta\|M_\beta^{-1/2}v\|^2
                         +\beta\|P_0x\|^2.
\]
By \eqref{E}, the form of \(E_{m,\beta}\) at the same vector is
exactly \(\beta\|U^*F_\beta\mathcal O_\beta v\|^2
=\beta\|P_0x\|^2\).
Subtracting proves \eqref{metriccomplement}.

Since \(G_\perp+\beta-M_\beta=CH_\beta^{-1}C^*\) is bounded,
their strictly positive form norms are equivalent for fixed
\(\beta\). This makes \(V_\beta\) and its inverse bounded and
proves the stated congruence.
The block inverse of the strictly positive gamma matrix yields
\[
 U^*(G+\beta)^{-1}U-H_\beta^{-1}
           =H_\beta^{-1}C^*M_\beta^{-1}CH_\beta^{-1}.
\]
Also
\[
 V_\beta^*V_\beta
     =I+M_\beta^{-1/2}CH_\beta^{-1}C^*M_\beta^{-1/2}.
\]
The nonzero spectra of the two finite-rank products give the
norm identity in \eqref{metricgrowth}.
For \(0<\beta\le1\), the eigenvalues of \(H_\beta\) lie between
\(\mu_m\) and \(\|H_0\|+1\). Theorem~\ref{main} and
\eqref{Enorm} now give the two-sided order in
\eqref{metricgrowth}.
\end{proof}

\section{Scope and relation to spectral positivity}

The positive pivot \(\beta H_\beta^{-1}\) has norm of order
\(\beta\) for a fixed frame. The subtracted correction
\eqref{E} has the larger order
\(\beta(\log(1/\beta))^{4m+5/2}\).
It nevertheless tends to zero. Thus a bound \(C_m\beta\) on this
particular correction cannot hold. A sign conclusion for
\eqref{complement} requires control of the retained complementary
block at the same scale or an exact cancellation within that block.
The correction alone is not a negative direction of the full Weil
form. Its magnitude also does not exclude a different normalization
or elimination strategy. Proposition~\ref{metric} gives such a
strategy explicitly: it removes the finite kernel block first,
while the same logarithmic growth appears in the change of gamma
metric. Neither normalization determines the sign of \(A_\perp\).

Recent work of Kim et al.\ \cite{Kim} treats numerical finite-interval
realizations of Suzuki's operator and its archimedean spectral law.
That setting and the present weighted whole-line resolvent have
different parameters and domains. The estimate here concerns
\(\beta\downarrow0\) for an exact fixed theta-jet span; it is not a
finite-interval eigenvalue computation.

The contribution proved here is the two-sided resolvent and
coupling order with a constructive range approximation. This
article makes no assertion of priority over all related weighted
resolvent estimates. It provides no uniform constants for a growing
frame, no matching pointwise spectral-mass lower bound, and no sign
for \(\mathcal C_{m,\beta}\). In particular it proves neither
unconditional positivity of all Li coefficients nor RH.

\begin{thebibliography}{9}
\bibitem{BL}
E.~Bombieri and J.~C.~Lagarias,
Complements to Li's criterion for the Riemann hypothesis,
J. Number Theory \textbf{77} (1999), 274--287.
\href{https://websites.umich.edu/~lagarias/doc/bombieri.pdf}
{Author manuscript}.
\bibitem{Romik}
D.~Romik,
Orthogonal polynomial expansions for the Riemann xi function,
\href{https://arxiv.org/abs/1902.06330v3}{arXiv:1902.06330v3}.
\bibitem{SuzukiLi}
M.~Suzuki,
Li coefficients as norms of functions in a model space.
\href{https://arxiv.org/abs/2301.05779v2}{arXiv:2301.05779v2} (2023).
\bibitem{SuzukiScrew}
M.~Suzuki,
Aspects of the screw function corresponding to the Riemann zeta-function,
J. London Math. Soc. \textbf{108} (2023), 1448--1487.
\href{https://arxiv.org/abs/2206.03682v4}{arXiv:2206.03682v4}.
\bibitem{SuzukiWeil}
M.~Suzuki,
Weil's quadratic form via the screw function.
\href{https://arxiv.org/abs/2606.09096v3}{arXiv:2606.09096v3} (2026).
\bibitem{Kim}
T.~Kim, Y.~Hong, M.~Kim, S.~Choi, J.~Jang, and M.~Kim,
A numerical realization of Suzuki's Weil-quadratic-form operator:
the archimedean spectral law, its universality, and an operator form
of Weil's positivity criterion.
\href{https://arxiv.org/abs/2607.24830v3}{arXiv:2607.24830v3} (2026).
\bibitem{Feshbach}
G.~Dusson, I.~M.~Sigal, and B.~Stamm,
The Feshbach--Schur map and perturbation theory.
\href{https://arxiv.org/abs/2105.02058v1}{arXiv:2105.02058v1} (2021).
\end{thebibliography}
\end{document}
